\documentclass[10pt,a4paper]{amsart}
\usepackage[margin=19mm]{geometry}
\usepackage{amsmath,amssymb,mathtools}
\usepackage[scr=rsfs,cal=euler]{mathalfa}
\usepackage{xfrac}
\usepackage[unicode,hidelinks,bookmarks=false]{hyperref}
\newcommand{\ie}{i.e.\ }
\newcommand{\cf}{cf.\ }
\newcommand{\resp}{respectively\ }
\newcommand{\Subsection}{\subsection}
\providecommand{\nobreakdash}{\nobreak\hbox{-}}
\setlength{\emergencystretch}{2em}
\multlinegap0pt
\usepackage{amssymb}
\usepackage{csquotes}
\usepackage[shortlabels]{enumitem}
\usepackage{mathtools}
\usepackage{tikz}
\usepackage{xcolor}
\newtheorem{theorem}{Theorem}
\newtheorem{lemma}{Lemma}
\newtheorem{claim}{Claim}
\usepackage{cleveref}
\crefname{theorem}{Theorem}{Theorems}
\crefname{lemma}{Lemma}{Lemmas}
\crefname{claim}{Claim}{Claims}
\title{Sufficient conditions for the bipartite rigidity, symmetric completability and hyperconnectivity of graphs}
\author{D\'aniel Garamv\"olgyi, Bill Jackson, Tibor Jord\'an, Soma Vill\'anyi}
\date{Source: arXiv:2511.00298v2 (CC BY 4.0)}
\begin{document}
\maketitle

\noindent\emph{Source and licence.} Sufficient conditions for the bipartite rigidity, symmetric completability and hyperconnectivity of graphs. Version arXiv:2511.00298v2 (CC BY 4.0), distributed by arXiv under the Creative Commons Attribution 4.0 International licence, http://creativecommons.org/licenses/by/4.0/, as recorded in the arXiv licence metadata for this exact version and shown on the abstract page https://arxiv.org/abs/2511.00298v2. Full licence notice: \texttt{licenses/arxiv-2511.00298-CC-BY-4.0.txt}.\par\smallskip

\noindent\textit{Editorial scope.} Counted unit: the article's theorem lem:tau\_bound2 (source0.tex lines 767--769). The article's complete original proof of it is delivered with it (source0.tex lines 771--818, one \texttt{proof} environment of 48 lines). No mathematics is written, repaired or re-derived: the statement and the proof are the article's own lines, in the article's own notation, and no retained line is substituted or re-flowed.  Retained uncounted as the source-local context the proof needs: lines 654--657; lines 729--735; lines 752--754.  Nothing was omitted from any retained span, so the package carries no omission record.  No retained line contains a citation, so the wrapper carries no bibliography and no bibliography entry.  No retained line contains a figure environment, an image-inclusion command or drawing code, so no image is delivered and the package declares no asset dependency.  Editorial formatting only, in addition: an AMS \texttt{amsart} preamble that restates the article's own declarations and macros used by the retained lines, namely the environments \texttt{theorem}, \texttt{lemma}, \texttt{claim} on separate counters, together with the standard packages those lines need.  The retained environments print in this document as Theorem 1, Lemma 1, Claim 1 in order of appearance, whereas the article prints the same items with the numbering of its own full document.  The complete native source of the article as distributed in the arXiv e-print for this exact version is bundled unchanged as \texttt{sources/188236/source0.tex}.
\begin{theorem}\label{lem:tau_bound}
    Let $G=(V,E)$ be a critically $k$-biconnected bipartite graph.
    Then $\tau(G)\geq \frac{|V|}{2k^2}$.
\end{theorem}

\begin{lemma}\label{cl}
    Let $G$, $\hat X$, $f$ and $G^f_X$ be as above. Then $G^f_X$ has $|\hat X|$ edges. In addition,
$E\cap E^f_X=\varnothing$
    and
    the multiplicity of each edge
$uv$ in $G^f_X$ is at most $k$.
\end{lemma}

\begin{claim}\label{cl2}
    For each $uv\in F$, we have $\kappa(u,v;G^+)\leq 2k-1$.
\end{claim}

\begin{theorem}\label{lem:tau_bound2}
    Let $G=(V,E)$ be a critically $k$-connected graph. Then $\tau(G)\geq \frac{|V|}{k+1}.$
\end{theorem}

\begin{proof}
We may assume that $k\geq 2$. First suppose that $|V|\leq 3k-2$.
Then $\frac{|V|}{k+1}<3$, so the theorem follows unless $k\leq \tau(G)\leq 2$.
Moreover, $k=\tau(G)$ holds only if $G$ contains $K_{k,|V|-k}$ as a spanning subgraph. Now the criticality of $G$ and $k=2$ imply that $G$ is 
either $K_3$ or $C_4$, for which the statement is clear.

Thus we may assume that $|V|\geq 3k-1$. The rest of the proof and 
our notation is similar to
that of \cref{lem:tau_bound}, except that 
we shall call a separator $S$ of $G$ 
{\itshape essential} if $|S|=k$ holds. 
Since $G$ is critically $k$-connected, for each vertex
$x\in V$ there exists an essential separator  $S$ with $x\in S$. 
Let $T\subseteq V$ be a smallest vertex cover and 
put $X=V-T$.

Next we define a pairing $f:X\to V\times V$ for $X$.
We choose a minimal sequence of essential separators $S_1,S_2,\dots,S_r$ such that
$X\subseteq \bigcup_1^r S$, and for each $j$, from $j=1$ to $j=r$,
we define $f(x)$ for
each $x\in X\cap (S_j-\bigcup_{i=1}^{j-1} S_i)$ sequentially as follows.
Fix $S_j$.
Since $S_j$
is a $k$-separator and $|V|\geq 3k-1$, there is a component $C$ of $G-S_j$
such that $|V(G-S_j-V(C))|\geq k.$ Let $D=V-S_j-V(C)$.
Let $G'$ be obtained from $G$ by
adding a new vertex $p$ and $k$ edges from $p$ to different vertices of $D$.
Let $q\in V(C)$. Now $G'$ is $k$-connected, and hence there exist $k$
internally disjoint paths from $p$ to $q$ in $G'$ by Menger's theorem.
Each vertex $x\in S_j-\bigcup_{i=1}^{j-1} S_i$ belongs to a subpath of length two in this collection
of $k$ paths, with end-vertices $u\in V(C)$ and $v\in D$.
Then we let $f(x)=(u,v)$.

Thus $f$ defines a graph $G_X^f=(V,E_X^f)$.
Since $T$ is a vertex cover of $G$ and $T\cap X=\emptyset$, each edge of $G_X^f$ is induced by $T$.
The key observation is that $G_X^f$ is simple.
The construction of
the pairing implies that the pairs defined for a fixed $S_j$ are
pairwise different. 
Suppose that for $y\in S_j$ we have $f(y)=(u,v)=f(x)$ for some
$x\in \bigcup_1^{j-1} S_i$, where $f(x)$ was defined earlier for some
$S_i$. The vertex $y$ is connected to both $u$ and $v$, and $u,v$ are separated by $S_i$,
which implies that $y\in S_i$, a contradiction. 

We may now use the argument in \cref{cl} and  \cref{cl2} to deduce that $G^+=G[T]\cup E_X^f$ is simple and,
for each $uv\in E_X^f$, we have $\kappa(u,v;G^+)\leq k$.
Then a similar count to that in the proof of \cref{lem:tau_bound} gives $\tau(G)\geq \frac{|V|}{k+1}$.
\end{proof}

\end{document}
